% Work/MAT -- nova hodina 37: slovni a technicke ulohy na mocniny a odmocniny.
% Po reorganizaci: 30 pravidla odmocnin, 31 racionalni exponent,
% 32 castecne odmocnovani, 33 usmernovani, 34 sdruzene vyrazy,
% 35 slozene vyrazy, 36 vedecky a technicky zapis, 37 aplikace.
% Ucebnice: Mocniny a Odmocniny.tex (prakticke motivace, vedecky zapis).
% Inspirace: Realisticky.cz 010608-010610, 020712-020715.
% Zadani jsou puvodni didakticke modely, ne realna mereni nebo citace.
\documentclass[11pt,a4paper]{article}
% Spolecny styl pro ucitelske pripravy MAT (Work/MAT)
% Verze 1.1, 2026-10-23
\NeedsTeXFormat{LaTeX2e}
\ProvidesPackage{priprava}[2026/10/23 v1.1 Sablona ucitelskych priprav MAT]

% Zakladni typografie a cestina
\RequirePackage[T1]{fontenc}
\RequirePackage[utf8]{inputenc}
\RequirePackage[czech]{babel}
\RequirePackage{lmodern}
\RequirePackage{microtype}
\RequirePackage{amsmath,amssymb}
\RequirePackage{xcolor}
\RequirePackage{enumitem}
\RequirePackage[most]{tcolorbox}
\RequirePackage{tabularx,booktabs,array}
\RequirePackage{fancyhdr}
\RequirePackage{needspace}
\RequirePackage{geometry}
\RequirePackage[hidelinks]{hyperref}

\geometry{a4paper,top=17mm,bottom=18mm,left=23mm,right=23mm}
\definecolor{notegray}{gray}{0.39}
\setlength{\parindent}{0pt}
\setlength{\parskip}{2.5pt plus 1pt minus 0.5pt}
\setlist[itemize]{topsep=2pt,itemsep=1pt,parsep=0pt,leftmargin=1.45em}
\setlist[enumerate]{topsep=2pt,itemsep=2pt,parsep=0pt,leftmargin=1.65em}
\setcounter{secnumdepth}{0}
\renewcommand{\arraystretch}{1.12}

% Zahlavi: vlevo rocnik | globalni tema, vpravo cislo: nazev hodiny
% Zapati zustava zcela prazdne.
\pagestyle{fancy}
\fancyhf{}
\renewcommand{\headrulewidth}{0.35pt}
\renewcommand{\footrulewidth}{0pt}
\setlength{\headheight}{14pt}
\newcommand{\nastavPripravu}[4]{%
  \fancyhead[L]{\small #1\enspace\textbar\enspace #2}%
  \fancyhead[R]{\small #3: #4}%
}

% Minimalisticke cernobile boxy. Kazda definice zvlast.
\tcbset{commonlesson/.style={enhanced,breakable=false,colback=white,colframe=black,
  colbacktitle=white,coltitle=black,boxrule=0.45pt,
  arc=0.8mm,outer arc=0.8mm,left=3mm,right=3mm,top=1.6mm,bottom=1.6mm,
  toptitle=1mm,bottomtitle=0.7mm,fonttitle=\bfseries,
  attach boxed title to top left={xshift=2.5mm,yshift=-1.15mm},
  boxed title style={colback=white,colframe=white,boxrule=0pt,left=1mm,right=1mm}}}
\newtcolorbox{zadaniBox}[1][]{commonlesson,title={Zadání},#1}
\newtcolorbox{prikladBox}[1][]{commonlesson,title={Příklad},#1}
\newtcolorbox{definiceBox}[1][]{commonlesson,title={Definice},#1}

% Spolecny box pro zadani a jeho reseni, bez mezery mezi castmi.
% Pouziti:
%   \begin{prikladReseniBox}
%     Zadani ulohy.
%     \tcblower
%     \begin{align*} ... \end{align*}
%   \end{prikladReseniBox}
% Spodni cast ma jemny sedy podklad a delici linku; lze v ni sazet
% vice radku matematiky. Obe casti maji jen jeden spolecny ramecek.
\tcbset{solutionlesson/.style={
  commonlesson,
  enhanced jigsaw,
  breakable,
  colbacklower=black!3,
  segmentation style={solid,draw=black!25,line width=0.35pt},
  fontlower=\small\color{notegray},
  before lower={\noindent\textbf{Řešení:}\enspace},
  leftlower=3mm,rightlower=3mm
}}
\newtcolorbox{prikladReseniBox}[1][]{solutionlesson,title={Příklad},#1}
\newtcolorbox{zadaniReseniBox}[1][]{solutionlesson,title={Zadání},#1}

% Pomucky pro autora pripravy: sedive poznamky; \reseni pro starsi pripravy
\newcommand{\poznamka}[1]{\par\begingroup\small\color{notegray}\itshape
  \textnormal{\textbf{Poznámka pro učitele.}}\enspace #1\par\endgroup}
\newcommand{\reseni}[1]{\par\begingroup\small\color{notegray}
  \textbf{Očekávané řešení:}\enspace #1\par\endgroup}
\newcommand{\blok}[1]{\needspace{5\baselineskip}\subsection*{#1}}
\newcommand{\vyrok}[1]{\uv{#1}}

\usepackage{tikz}
\nastavPripravu{1. ročník}{Mocniny a odmocniny}{37}{Slovní a technické úlohy}
\begin{document}
\poznamka{\textbf{45minutové jádro:} 4 min převod textu na schéma;
9 min Příklad 1; 8 min Samostatná práce 1 a), b);
8 min Příklad 2; 9 min Samostatná práce 2 a), b);
4 min Příklad 3 (vědecký zápis a jednotky); 3 min shrnutí.
Ostatní úlohy, příklady z informatiky a slovní domácí úkoly tvoří zásobník.
\textbf{Důležité:} číselné údaje jsou modelové, nikoliv naměřené nebo aktuální ceny.}

\blok{Dnes se neptáme, který vzorec je napsaný nad příkladem}
Úkolem je z textu rozpoznat potřebnou matematiku a \textbf{zvolit správný postup}.
Výsledek bez jednotky nebo bez kontroly smyslu není hotová odpověď.

\begin{definiceBox}[title={Postup technické úlohy -- praktická strategie}]
\begin{enumerate}[label=\arabic*.]
\item \textbf{Rozuměj situaci:} načrtni obrázek, vypiš známé údaje.
\item \textbf{Sjednoť jednotky:} například $1\,\mathrm{m}=100\,\mathrm{cm}$,
$1\,\mathrm{m}^3=1000\,\mathrm{l}$.
\item \textbf{Vytvoř model:} mocnina, odmocnina, Pythagorova věta nebo vědecký zápis.
\item \textbf{Vypočítej a zkontroluj:} přesný výsledek před zaokrouhlením, správná jednotka.
\end{enumerate}
\end{definiceBox}

\blok{Geometrie: odmocnina umí vypočítat rozměr}
\begin{prikladReseniBox}[title={Příklad 1: úhlopříčka ochranného panelu}]
Na dílně se úhlopříčně vyztužuje obdélníkový rám o šířce
$2{,}4\,\mathrm m$ a výšce $1{,}8\,\mathrm m$. Jak dlouhá je samotná
výztuha? Při výrobě se počítá s \emph{desetiprocentní rezervou délky}.
Jak dlouhý kus materiálu máme připravit?
\tcblower
\begin{center}
\begin{tikzpicture}[x=1cm,y=1cm,scale=.8]
\draw[thick] (0,0) rectangle (5,3);
\draw[very thick] (0,0)--(5,3);
\node[below] at (2.5,0) {$2{,}4\,\mathrm m$};
\node[right] at (5,1.5) {$1{,}8\,\mathrm m$};
\node[above,rotate=31] at (2.7,1.6) {$d$};
\end{tikzpicture}
\end{center}
Pythagorova věta dává
\[
d=\sqrt{2{,}4^2+1{,}8^2}=\sqrt{5{,}76+3{,}24}=\sqrt9=3\,\mathrm m.
\]
S rezervou potřebujeme $3\cdot1{,}10=\boxed{3{,}3\,\mathrm m}$ materiálu.
Odhad: úhlopříčka musí být delší než obě strany, což $3\,\mathrm m$ splňuje.
\end{prikladReseniBox}

\begin{zadaniReseniBox}[title={Samostatná práce 1: rozměry a jednotky}]
\begin{enumerate}[label=\alph*)]
\item Čtvercový plech má obsah $242\,\mathrm{cm}^2$. Urči přesně délku
strany a obvod; délku strany pak přibližně na desetiny centimetru.
\item Krychlová nádoba má objem $27\,\mathrm{dm}^3$.
Urči délku hrany v decimetrech a centimetrech.
\item (Navíc) Příčná vzpěra spojuje dva body vzdálené vodorovně $12\,\mathrm{dm}$
a svisle $5\,\mathrm{dm}$. Jak dlouhá je přímá vzpěra?
\end{enumerate}
\tcblower
\begin{enumerate}[label=\alph*)]
\item $a=\sqrt{242}=11\sqrt2\,\mathrm{cm}\doteq15{,}6\,\mathrm{cm}$;
$o=44\sqrt2\,\mathrm{cm}\doteq62{,}2\,\mathrm{cm}$.
\item $a=\sqrt[3]{27}=3\,\mathrm{dm}=30\,\mathrm{cm}$.
\item $d=\sqrt{12^2+5^2}=\sqrt{169}=13\,\mathrm{dm}$.
\end{enumerate}
\end{zadaniReseniBox}

\blok{Technika: měřítko a řád velikosti}
\begin{prikladReseniBox}[title={Příklad 2: tenký drát a násobky tisíciny}]
Jmenovitý průměr drátku je $0{,}00035\,\mathrm m$.
Zapiš jej ve vědeckém i technickém tvaru a v milimetrech a mikrometrech.
\tcblower
\[
0{,}00035\,\mathrm m=3{,}5\cdot10^{-4}\,\mathrm m
=350\cdot10^{-6}\,\mathrm m
=0{,}35\,\mathrm{mm}=350\,\mathrm{\mu m}.
\]
Vědecký tvar má koeficient s absolutní hodnotou mezi $1$ a $10$;
v technickém tvaru je exponent násobkem tří.
\end{prikladReseniBox}

\begin{zadaniReseniBox}[title={Samostatná práce 2: rychlý odhad a výkon}]
\begin{enumerate}[label=\alph*)]
\item Elektromagnetický signál se šíří ve vakuu rychlostí přibližně
$3\cdot10^8\,\mathrm{m/s}$. Jakou vzdálenost urazí za $2{,}5\,\mathrm{ms}$?
Převeď na kilometry.
\item Topný rezistor má odpor $6\,\Omega$ a je na něm napětí $12\,\mathrm V$.
Při použití vztahu $P=U^2/R$ vypočti příkon ve wattech.
Kolik energie spotřebuje za dvě hodiny v \(\mathrm{kWh}\)?
\item (Navíc) Zaokrouhli $0{,}0000789\,\mathrm m$ na dvě platné číslice;
zapiš výsledek pomocí vhodné předpony jednotky.
\end{enumerate}
\tcblower
\begin{enumerate}[label=\alph*)]
\item $t=2{,}5\cdot10^{-3}\,\mathrm s$,
$s=3\cdot10^8\cdot2{,}5\cdot10^{-3}=7{,}5\cdot10^5\,\mathrm m=750\,\mathrm{km}$.
\item $P=12^2/6=24\,\mathrm W$;
$E=24\cdot2=48\,\mathrm{Wh}=0{,}048\,\mathrm{kWh}$.
\item $0{,}000079\,\mathrm m=79\,\mathrm{\mu m}$.
\end{enumerate}
\end{zadaniReseniBox}

\poznamka{Vzorec $P=U^2/R$ i přibližná rychlost světla jsou v zadání;
nejde o předpoklad znalosti celé fyziky. Požaduj kontrolu jednotek
ve výpočtu a sdělení, zda je výsledná hodnota realistická.}

\blok{Informatika: mocniny dvojky}
\begin{prikladReseniBox}[title={Příklad 3: záznamy ze senzorů},breakable=false]
Systém uloží data ze $128$ senzorů. Každý senzor vytvoří soubor o velikosti
$64\,\mathrm{KiB}$. Kolik je to dohromady bajtů a mebibajtů ($\mathrm{MiB}$)?
Uvažuj $1\,\mathrm{KiB}=2^{10}\,\mathrm B$ a $1\,\mathrm{MiB}=2^{20}\,\mathrm B$.
\tcblower
$128=2^7$, $64=2^6$; celkem
\[
2^7\cdot2^6\cdot2^{10}=2^{23}=8\,388\,608\,\mathrm B
=\boxed{8\,\mathrm{MiB}}.
\]
Nezaměňuj $\mathrm{MB}=10^6\,\mathrm B$ a $\mathrm{MiB}=2^{20}\,\mathrm B$.
\end{prikladReseniBox}

\begin{zadaniReseniBox}[title={Samostatná práce 3: smíšené technické úlohy (rezerva)}]
\begin{enumerate}[label=\alph*)]
\item Obrázek má $1920\times1080$ pixelů; na jeden pixel jsou použity
$3$ bajty. Kolik bajtů potřebuje nekomprimovaný obraz bez hlavičky?
Zaokrouhli velikost na dvě desetinná místa v MB.
\item Nádrž ve tvaru kvádru měří $1{,}2\,\mathrm m\times0{,}5\,\mathrm m\times0{,}8\,\mathrm m$.
Kolik litrů vody se do ní teoreticky vejde, jestliže ji naplníme po okraj?
\item Čtvercová podložka má obsah $50\,\mathrm{cm}^2$. Napiš její
stranu přesně pomocí odmocniny a obvod ve stejném tvaru.
\end{enumerate}
\tcblower
\begin{enumerate}[label=\alph*)]
\item $1920\cdot1080\cdot3=6\,220\,800\,\mathrm B\doteq6{,}22\,\mathrm{MB}$.
\item $V=1{,}2\cdot0{,}5\cdot0{,}8=0{,}48\,\mathrm{m}^3=480\,\mathrm l$.
\item $a=\sqrt{50}=5\sqrt2\,\mathrm{cm}$,
$o=4a=20\sqrt2\,\mathrm{cm}$.
\end{enumerate}
\end{zadaniReseniBox}

\blok{Rozšiřující příklad 1: závislost na druhé mocnině}
\begin{prikladReseniBox}[title={Rozšiřující příklad 1: intenzita při změně vzdálenosti},breakable=false]
V idealizovaném modelu bodového zdroje je intenzita záření nepřímo
úměrná druhé mocnině vzdálenosti: $I=k/r^2$, kde $k>0$.
Kolikrát se změní intenzita, jestliže vzdálenost zvětšíme ze
$2\,\mathrm m$ na $6\,\mathrm m$?
\tcblower
$\dfrac{I(6)}{I(2)}=\dfrac{k/6^2}{k/2^2}=
\dfrac{2^2}{6^2}=\dfrac19$.
Intenzita je \textbf{devětkrát menší}, nikoli třikrát.
Model nepopisuje všechny skutečné zdroje nebo prostředí.
\end{prikladReseniBox}

\blok{Domácí úkol 1: tři různé situace}
\begin{zadaniReseniBox}[title={Domácí úkol 1: nezapomeň na jednotky},breakable=false]
\begin{enumerate}[label=\alph*)]
\item Krychlová nádrž má objem $216\,\mathrm l$. Jaká je délka hrany v centimetrech?
\item Čtvercový pozemek má obsah $72\,\mathrm{m}^2$. Vypočti délku strany a
obvod přesně pomocí odmocniny.
\item Procesor provede $2{,}4\cdot10^9$ operací za sekundu.
Kolik operací podle tohoto zjednodušeného modelu vykoná za $5\,\mathrm{ms}$?
\end{enumerate}
\tcblower
\begin{enumerate}[label=\alph*)]
\item $216\,\mathrm l=216\,\mathrm{dm}^3$, tedy $a=\sqrt[3]{216}=6\,\mathrm{dm}=60\,\mathrm{cm}$.
\item $a=\sqrt{72}=6\sqrt2\,\mathrm m$;
$o=24\sqrt2\,\mathrm m$.
\item $2{,}4\cdot10^9\cdot5\cdot10^{-3}=12\cdot10^6=1{,}2\cdot10^7$ operací.
\end{enumerate}
\end{zadaniReseniBox}

\poznamka{Závěrečná otázka: Kdy je vhodnější nechat výsledek jako
$6\sqrt2\,\mathrm m$ a kdy jej zaokrouhlit na centimetry?
Důležité je rozlišovat přesný matematický výsledek a technický údaj
s danou přesností měření. Příští hodina může být opakování před testem.}
\end{document}
